\begin{center}
  \large\textbf{ABSTRAK}
\end{center}

\addcontentsline{toc}{chapter}{ABSTRAK}

\begin{center}
  \textbf{\tatitle{}}
\end{center}

\vspace{1ex}

\begingroup
\setlength{\tabcolsep}{0pt}

\noindent
\begin{tabularx}{\textwidth}{l >{\centering}m{2em} X}
  Nama Mahasiswa / NRP & : & \name{} / \nrp{} \\
  Departemen           & : & \department{} \facultyshort{} - ITS \\
  Dosen Pembimbing     & : & 1. \advisor{} \\
                       &   & 2. \coadvisor{} \\
\end{tabularx}
\endgroup

\vspace{2ex}

\noindent\textbf{Abstrak}

% Ubah paragraf berikut dengan abstrak dari proposal tugas akhir
Pada penelitian ini diajukan \lipsum[1]

\vspace{2ex}

\noindent\textbf{Kata kunci:} Edge AI, \emph{Speech-to-Text}, SLM, SOAP, Rekam Medis.

